% === START TEMPLATE WITH CONFIGURABLE WATERMARK POSITION ===
\documentclass[10pt]{article}
\usepackage[utf8]{inputenc}
\usepackage[left=50pt,right=50pt,top=100pt,bottom=100pt]{geometry}
\usepackage[T1]{fontenc}
\usepackage{lmodern}
\usepackage[hidelinks]{hyperref}
\usepackage{graphicx}
\usepackage[table]{xcolor}
\usepackage{transparent}      % for opacity
\usepackage{tabularx}
\usepackage{stackengine}
\usepackage{mwe}             % example-image, example-image-a, etc.
\usepackage{tocloft}
\usepackage[hypcap=false,font=small]{caption}
\usepackage{capt-of}         % for captionof outside floats
\usepackage{multicol}
\usepackage{ragged2e}
\usepackage{titlesec}
\usepackage{fancyhdr}
\usepackage{etoolbox}
\usepackage{eso-pic}         % for background images

% === User-Configurable Variables ===
\newcommand{\ProductName}{PRODUCT NAME HERE}
\newcommand{\ProductNameShort}{SHORTNAME}
\newcommand{\WebsiteLinkText}{www.example.com}
\newcommand{\ReleaseDate}{\today}
\newcommand{\Revision}{Rev.\ X.Y}

% === Watermark Variables ===
\newcommand{\WatermarkPath}{example-image}   % built-in placeholder
\newcommand{\WatermarkScale}{0}              % scale factor (0 to disable watermark)
\newcommand{\WatermarkAngle}{45}             % rotation angle
\newcommand{\WatermarkOpacity}{0.1}          % 0 (transparent) to 1 (opaque)
\newcommand{\WatermarkPosX}{0.1\paperwidth}  % horizontal starting point
\newcommand{\WatermarkPosY}{0.1\paperheight} % vertical starting point

% Place watermark on every page using configurable origin
\AddToShipoutPictureBG{%
	\AtPageLowerLeft{%
		\raisebox{\WatermarkPosY}[0pt][0pt]{%
			\makebox[\WatermarkPosX][l]{%
				\transparent{\WatermarkOpacity}%
				\includegraphics[scale=\WatermarkScale,angle=\WatermarkAngle]{\WatermarkPath}%
			}%
		}%2
	}%
}

% === TOC configuration ===
\renewcommand{\contentsname}{Table of Contents}
\setcounter{tocdepth}{2}
\renewcommand{\cftdotsep}{1.5}
\renewcommand{\cftsecleader}{\cftdotfill{\cftdotsep}}
\renewcommand{\cftsubsecleader}{\cftdotfill{\cftdotsep}}
\setlength{\cftbeforesecskip}{0.5em}
\setlength{\cftbeforesubsecskip}{0.5em}

% === Sans-serif font ===
\renewcommand{\familydefault}{\sfdefault}

% === Header/Footer styling ===
\pagestyle{fancy}
\fancyhf{}
\renewcommand{\headrulewidth}{0.4pt}
\renewcommand{\footrulewidth}{0.4pt}
\setlength{\headheight}{60pt}
\setlength{\headsep}{10pt}
\setlength{\footskip}{50pt}
\lhead{\Large\textbf{\ProductNameShort}}
\rhead{%
	\includegraphics[height=0.6cm]{example-image-a}\\[-0.5em]
	{\scriptsize\textcolor{gray!50}{\WebsiteLinkText}}%
}
\fancyfoot[L]{\footnotesize\textcolor{gray!50}{\ReleaseDate}}
\fancyfoot[C]{\footnotesize\textcolor{gray!50}{\Revision}}
\fancyfoot[R]{\footnotesize\textcolor{gray!50}{\thepage}}

\fancypagestyle{plain}{%
	\fancyhf{}%
	\lhead{\Large\textbf{\ProductNameShort}}%
	\rhead{%
		\includegraphics[height=0.6cm]{example-image-a}\\[-0.5em]
		{\scriptsize\textcolor{gray!50}{\WebsiteLinkText}}%
	}%
	\fancyfoot[L]{\footnotesize\textcolor{gray!50}{\ReleaseDate}}%
	\fancyfoot[C]{\footnotesize\textcolor{gray!50}{\Revision}}%
	\fancyfoot[R]{\footnotesize\textcolor{gray!50}{\thepage}}%
	\renewcommand{\headrulewidth}{0.4pt}%
	\renewcommand{\footrulewidth}{0.4pt}%
}

% === Section/Subsection numbering & formatting ===
\setcounter{secnumdepth}{3}
\renewcommand{\thesubsection}{\arabic{section}.\arabic{subsection}}
\renewcommand{\thesubsubsection}{\arabic{section}.\arabic{subsection}.\arabic{subsubsection}}
\titleformat{\subsection}[hang]
{\large\bfseries}{\thesubsection}{1em}{}
[\color{gray!30}\rule{\textwidth}{0.4pt}]
\titlespacing*{\subsection}{0pt}{1.5\baselineskip}{1\baselineskip}

\newcommand{\sectionbreak}{\clearpage}

\begin{document}
	\flushbottom
	
	%% Title Page
	\begin{center}
		{\LARGE\bfseries \ProductName\par}
		\vspace{1ex}
		{\normalsize Datasheet Template\par}
	\end{center}
	\vspace{2\baselineskip}
	
	%% Front page content
	\noindent
	\begin{minipage}[t]{0.45\textwidth}
		\section{Features}\label{sec:features}
		\vspace{\baselineskip}
		\begin{itemize}
			\item -- FEATURE A
			\item -- FEATURE B
			\item -- FEATURE C
			\item -- FEATURE D
		\end{itemize}
	\end{minipage}
	\hfill
	\begin{minipage}[t]{0.5\textwidth}
		\section{Description}\label{sec:desc}
		\vspace{\baselineskip}
		This template provides a structured outline for a generic datasheet,
		including configurable headers, footers, and a background watermark.
		\par\medskip
		\begin{center}
			\includegraphics[width=5cm]{example-image}
			\captionof{figure}{Module Overview}\label{fig:component}
		\end{center}
	\end{minipage}
	
	%% TOC + Revision History (same page)
	\section*{\contentsname}
	\begin{multicols}{2}
		{\let\oldcontentsname\contentsname
			\renewcommand{\contentsname}{}
			\tableofcontents
			\let\contentsname\oldcontentsname}
	\end{multicols}
	
	\vspace{1em}
	\noindent\textbf{Revision History}\par
	\begin{tabularx}{\textwidth}{|l|X|}
		\hline
		\rowcolor{gray!25}\textbf{Date} & \textbf{Description} \\ \hline
		\ReleaseDate & Initial release of template \\ \hline
	\end{tabularx}
	\captionof{table}{Revision History}\label{tab:rev}
	
	% Main sections each on new page:
	\sectionbreak
	
	\section{Device Overview}\label{sec:overview}
	\subsection{Part Number Options}\label{sec:overview.1}
	\begin{tabularx}{\textwidth}{|X|X|X|}
		\hline
		\rowcolor{gray!25}\textbf{PART NUMBER} & \textbf{PACKAGE} & \textbf{DESCRIPTION} \\ \hline
		PART-0001 & GENERIC SIZE & Generic description \\ \hline
	\end{tabularx}
	\captionof{table}{Part Number Options}\label{tab:overview1}
	
	\subsection{Pin Configuration and Functions}\label{sec:overview.2}
	\begin{center}
		\includegraphics[width=0.6\textwidth]{example-image-a}
		\captionof{figure}{Pinout Diagram}\label{fig:pinout}
	\end{center}
	\begin{itemize}
		\item PIN1 -- Generic function
		\item PIN2 -- Generic function
		\item PIN3 -- Generic function
		\item PIN4 -- Generic function
	\end{itemize}
	
	\sectionbreak
	
	\section{Specifications}\label{sec:specs}
	\subsection{Absolute Maximum Ratings}\label{sec:specs.1}
	\begin{tabularx}{\textwidth}{|X|X|X|}
		\hline
		\rowcolor{gray!25}\textbf{PARAMETER} & \textbf{MAX RATING} & \textbf{UNIT} \\ \hline
		PARAM A & VALUE RANGE & UNIT \\ \hline
		PARAM B & VALUE RANGE & UNIT \\ \hline
	\end{tabularx}
	\captionof{table}{Absolute Maximum Ratings}\label{tab:amax}
	
	\subsection{Recommended Operating Conditions}\label{sec:specs.2}
	\begin{tabularx}{\textwidth}{|X|X|X|}
		\hline
		\rowcolor{gray!25}\textbf{PARAMETER} & \textbf{TYPICAL} & \textbf{UNIT} \\ \hline
		PARAM C & TYPICAL & UNIT \\ \hline
		PARAM D & TYPICAL & UNIT \\ \hline
	\end{tabularx}
	\captionof{table}{Recommended Operating Conditions}\label{tab:opcond}
	
	\subsection{Electrical Characteristics}\label{sec:specs.3}
	\begin{tabularx}{\textwidth}{|X|X|X|}
		\hline
		\rowcolor{gray!25}\textbf{PARAMETER} & \textbf{TYPICAL} & \textbf{UNIT} \\ \hline
		PARAM E & TYPICAL & UNIT \\ \hline
		PARAM F & TYPICAL & UNIT \\ \hline
	\end{tabularx}
	\captionof{table}{Electrical Characteristics}\label{tab:elec}
	
	\sectionbreak
	
	\section{Typical Characteristics}\label{sec:typ}
	\begin{center}
		\includegraphics[width=\textwidth]{example-image-b}
		\captionof{figure}{Characteristic Chart 1}\label{fig:chart1}
		\medskip
		
		\includegraphics[width=\textwidth]{example-image-c}
		\captionof{figure}{Characteristic Chart 2}\label{fig:chart2}
	\end{center}
	
	\sectionbreak
	
	\section{Functional Block Diagram}\label{sec:block}
	\begin{center}
		\includegraphics[width=0.8\textwidth]{example-image}
		\captionof{figure}{Block Diagram}\label{fig:block}
	\end{center}
	
	\sectionbreak
	
	\section{Feature Description}\label{sec:features_desc}
	\subsection{Feature A}\label{sec:features_desc.1}
	Generic feature description.
	\subsection{Feature B}\label{sec:features_desc.2}
	Generic feature description.
	\subsection{Feature C}\label{sec:features_desc.3}
	Generic feature description.
	\subsection{Feature D}\label{sec:features_desc.4}
	Generic feature description.
	
	\sectionbreak
	
	\section{Application Information}\label{sec:app}
	\subsection{Typical Application}\label{sec:app.1}
	\begin{center}
		\includegraphics[width=\textwidth]{example-image-a}
		\captionof{figure}{Application Diagram}\label{fig:app}
	\end{center}
	\subsection{Design Considerations}\label{sec:app.2}
	\begin{itemize}
		\item -- Generic consideration A
		\item -- Generic consideration B
	\end{itemize}
	
	\sectionbreak
	
	\section{Power Supply Recommendations}\label{sec:ps}
	Generic power supply recommendation text.
	
	\sectionbreak
	
	\section{Layout Guidelines}\label{sec:layout}
	\subsection{PCB Recommendations}\label{sec:layout.1}
	Generic PCB layout guidance.
	\subsection{Thermal Considerations}\label{sec:layout.2}
	Generic thermal management guidance.
	
	\sectionbreak
	
	\section{Package Information}\label{sec:package}
	\subsection{Mechanical Dimensions}\label{sec:package.1}
	\begin{center}
		\includegraphics[width=0.8\textwidth]{example-image-b}
		\captionof{figure}{Mechanical Drawing}\label{fig:mech}
	\end{center}
	\subsection{Land Pattern Recommendations}\label{sec:package.2}
	\begin{itemize}
		\item Generic land pattern recommendation A
		\item Generic land pattern recommendation B
	\end{itemize}
	
	\sectionbreak
	
	\section{Ordering Information}\label{sec:order}
	\begin{tabularx}{\textwidth}{|X|X|X|}
		\hline
		\rowcolor{gray!25}\textbf{PART NUMBER} & \textbf{ORDER CODE} & \textbf{DESCRIPTION} \\ \hline
		PART-0001 & ORDER-1234 & Generic ordering info \\ \hline
	\end{tabularx}
	\captionof{table}{Ordering Information}\label{tab:order}
	
\end{document}
